\documentclass[12pt, a4paper]{article}
\usepackage[utf8]{inputenc}
\usepackage[T1]{fontenc}
\usepackage[vietnamese]{babel}
\usepackage{geometry}
\geometry{a4paper, margin=2cm}
\usepackage{longtable}
\usepackage{array}
\usepackage{booktabs}

\title{\textbf{BrewLite -- Kế hoạch Sprint (Sprint Backlog)}}
\author{Nhóm BrewLite}
\date{Học kỳ Thu 2026}

\begin{document}
\maketitle

\section{Tổng quan}
Tài liệu này là \textbf{Kế hoạch Sprint} (Sprint Backlog) cho dự án BrewLite, dựa trên các yêu cầu của \textbf{Bài tập lớn 2 -- BrewLite -- Phiên bản 1.0} (Môn Công nghệ phần mềm, học kỳ I -- năm học 2026-2027).

Dự án xây dựng hệ thống đặt món cà phê (không tiền mặt, ưu tiên thiết bị di động/PWA) với:
frontend sử dụng Next.js, backend sử dụng NestJS, cơ sở dữ liệu PostgreSQL, cổng thanh toán giả lập (mock payment gateway).

Luồng nghiệp vụ khả thi tối thiểu (MVP):
\begin{center}
\texttt{Xem menu -> Chọn tùy chọn -> Giỏ hàng -> Đặt đơn -> Thanh toán -> Theo dõi đơn}
\end{center}

\section{Sprint 1}
\label{sec:sprint1}

\subsection{Mục tiêu Sprint}
\textbf{Nền tảng \& Hiển thị thực đơn} -- Xây dựng nền tảng cơ bản và hiển thị thực đơn.

\subsection{Các công việc được giao cho Sprint 1}
\begin{itemize}
  \item Nhiệm vụ 1 -- Khảo sát và đặc tả yêu cầu (ngữ cảnh/phạm vi, tài liệu SRS, sơ đồ use-case).
  \item Nhiệm vụ 2 -- Mô hình hóa hệ thống (sơ đồ hoạt động, sơ đồ tuần tự, sơ đồ lớp).
  \item Nhiệm vụ 3 -- Thiết kế kỹ thuật (C4, quyết định kiến trúc ADR, sơ đồ triển khai).
  \item Nhiệm vụ 4 -- Triển khai Sprint 1 (khởi tạo kho mã nguồn, tài liệu, menu MVP, kiểm thử liên tục CI).
\end{itemize}

\subsection{Danh mục công việc và sản phẩm bàn giao trong Sprint Backlog}
\begin{enumerate}
  \item \textbf{Nhiệm vụ 1} -- Tài liệu yêu cầu:
  \begin{itemize}
    \item Ngữ cảnh và phạm vi (các bên liên quan, phạm vi trong/ngoài dự án, các giả định).
    \item Tài liệu SRS (FR-01..FR-16, NFR-01..NFR-08, mô hình MoSCoW, tiêu chí nghiệm thu).
    \item Use-case UC-07 -- Đặt và thanh toán đơn hàng.
  \end{itemize}
  \item \textbf{Nhiệm vụ 2} -- Tài liệu mô hình hóa:
  \begin{itemize}
    \item Sơ đồ hoạt động (Activity diagram) cho toàn bộ luồng nghiệp vụ.
    \item Sơ đồ tuần tự (Sequence diagram) cho UC-07 (JWT, giao dịch, tính idempotency, thanh toán).
    \item Sơ đồ lớp (Class diagram: User, Product, Order, OrderItem, Payment, Promotion, LoyaltyTransaction).
  \end{itemize}
  \item \textbf{Nhiệm vụ 3} -- Tài liệu kiến trúc:
  \begin{itemize}
    \item Mô hình C4 (Context / Container / Component).
    \item Các tài liệu quyết định kiến trúc ADR 001--006.
    \item Sơ đồ triển khai cho môi trường sản xuất (Render).
  \end{itemize}
  \item \textbf{Nhiệm vụ 4} -- Kỹ thuật Sprint 1:
  \begin{itemize}
    \item Monorepo gồm apps/web + apps/api + tài liệu.
    \item NestJS: GET /products, GET /products/:id, module Kiểm tra hệ thống (Health).
    \item Giao diện menu Next.js (ưu tiên di động, trạng thái tải/trống/lỗi).
    \item Lược đồ Prisma, tập lệnh di chuyển cơ sở dữ liệu (migration) và dữ liệu mẫu (seed).
    \item Cổng chất lượng CI (Cài đặt -> Kiểm tra lỗi cú pháp -> Chạy Unit Test -> Xây dựng bản dựng).
  \end{itemize}
\end{enumerate}

\subsection{Tiêu chí nghiệm thu Sprint 1 (Tóm tắt)}
\begin{itemize}
  \item Hoàn thành đầy đủ các yêu cầu, sơ đồ, tài liệu ADR, v.v. lên kho mã nguồn.
  \item Mô hình hóa không có điểm mâu thuẫn với tài liệu SRS.
  \item Bản mẫu MVP chạy được luồng "xem menu".
\end{itemize}

\section{Sprint 2}
\label{sec:sprint2}

\subsection{Mục tiêu Sprint}
\textbf{Giỏ hàng, Đặt đơn \& Đăng nhập} / \textbf{Thanh toán, Bàn giao \& Nghiệp vụ nâng cao}.

\subsection{Các công việc được giao cho Sprint 2}
\begin{itemize}
  \item Nhiệm vụ 5 -- Hoàn thiện MVP (chi tiết sản phẩm, kích thước/topping, giỏ hàng, xác thực JWT, đặt đơn, thanh toán ví/thẻ giả lập, xác nhận/trạng thái/lịch sử, xử lý lỗi và kiểm tra dữ liệu).
  \item Nhiệm vụ 6 -- Tích hợp và kiểm thử (unit test, integration test, E2E, docker-compose với 3 dịch vụ, endpoint kiểm tra sức khỏe hệ thống).
  \item Nhiệm vụ 7 -- Triển khai liên tục CD \& Bàn giao (pipeline CD, di chuyển CSDL an toàn, triển khai, kiểm thử khói, bàn giao).
\end{itemize}

\subsection{Danh mục công việc và sản phẩm bàn giao trong Sprint Backlog}
\begin{enumerate}
  \item \textbf{Các tính năng MVP}: giỏ hàng (Zustand + localStorage), đặt hàng (POST /orders), thanh toán (POST /payments + cổng thanh toán giả lập, tính đảm bảo idempotency), xác nhận và lịch sử đơn hàng (GET /orders/me).
  \item \textbf{Kiểm thử}: unit test cho nghiệp vụ (máy trạng thái, tính giá, idempotency, khuyến mại, điểm tích lũy); integration test cho quy trình đặt đơn/thanh toán; kiểm thử E2E từ menu đến xác nhận.
  \item \textbf{CD}: Pipeline đẩy mã/gắn thẻ (Push/Tag) -> Kiểm tra lỗi \& Test -> Xây dựng -> Đẩy hình ảnh Docker lên Registry -> Triển khai -> Kiểm thử khói; chạy di chuyển CSDL an toàn; cơ chế khôi phục (rollback).
\end{enumerate}

\section{Sprint 3}
\label{sec:sprint3}

\subsection{Mục tiêu Sprint}
\textbf{Nghiệp vụ backend nâng cao \& Bàn giao hoàn chỉnh}.

\subsection{Các công việc được giao cho Sprint 3}
\begin{itemize}
  \item Nhiệm vụ 8 -- Thanh toán không tiền mặt (cổng thanh toán giả lập: Ví/Thẻ, cơ chế idempotency, xử lý lỗi \texttt{PAYMENT\_FAILED}).
  \item Nhiệm vụ 9 -- Xác nhận, lịch sử đơn hàng \& Bàn giao (docker-compose chính thức, file README hướng dẫn, bản demo E2E).
  \item Nhiệm vụ 10 -- Nghiệp vụ backend chuyên sâu (Máy trạng thái đơn hàng, idempotency, khóa tồn kho, khuyến mại/điểm thưởng).
\end{itemize}

\subsection{Danh mục công việc và sản phẩm bàn giao trong Sprint Backlog}
\begin{itemize}
  \item Logic nghiệp vụ backend: Máy trạng thái chuyển đổi \texttt{PENDING -> PAID -> PREPARING -> READY -> COMPLETED}, hoặc \texttt{PENDING -> PAYMENT\_FAILED -> PAID}, và \texttt{CANCELLED}.
  \item Ràng buộc khóa duy nhất \texttt{Idempotency-Key UNIQUE} cho bảng Thanh toán.
  \item Khóa lạc quan (\texttt{Product.version}) dùng cho quản lý kho hàng.
  \item Điểm tích lũy thành viên (\texttt{floor(total / 1000)}) kết hợp xử lý mã khuyến mại.
\end{itemize}

\section{Bảng tổng hợp Sprint}

\begin{longtable}{p{3cm} p{3.5cm} p{4.5cm}}
\toprule
\textbf{Sprint} & \textbf{Mục tiêu} & \textbf{Công việc / Trọng tâm} \\
\midrule
Sprint 1 & Nền tảng \& Hiển thị thực đơn & Nhiệm vụ 1, 2, 3, 4 (Yêu cầu, Mô hình hóa, Kỹ thuật, Menu MVP) \\
Sprint 2 & Giỏ hàng, Đặt đơn, Đăng nhập / Thanh toán, Bàn giao \& Nghiệp vụ nâng cao & Nhiệm vụ 5, 6, 7 (Hoàn thiện MVP, Tích hợp/Kiểm thử, CD/Triển khai/Bàn giao) \\
Sprint 3 & Nghiệp vụ backend \& Bàn giao & Nhiệm vụ 8, 9, 10 (Thanh toán, Xác nhận/Lịch sử/Bàn giao, Máy trạng thái/Idempotency/Tồn kho) \\
\bottomrule
\end{longtable}

\section{Tiêu chí hoàn thành (Definition of Done - tóm tắt)}
\begin{itemize}
  \item Mã nguồn chạy được, không lỗi build; tuân thủ toàn bộ các tiêu chuẩn của dự án.
  \item Có commit trên Git với mô tả rõ ràng, minh bạch.
  \item API có kiểm tra tính hợp lệ dữ liệu (validate); giao diện xử lý tốt trạng thái tải (loading) và lỗi (error).
  \item Có hướng dẫn chạy ứng dụng đầy đủ trong tệp README.
\end{itemize}

\end{document}
